%%=============================================================================
%% Samenvatting
%%=============================================================================

% TODO: De "abstract" of samenvatting is een kernachtige (~ 1 blz. voor een
% thesis) synthese van het document.
%
% Een goede abstract biedt een kernachtig antwoord op volgende vragen:
%
% 1. Waarover gaat de bachelorproef?
% 2. Waarom heb je er over geschreven?
% 3. Hoe heb je het onderzoek uitgevoerd?
% 4. Wat waren de resultaten? Wat blijkt uit je onderzoek?
% 5. Wat betekenen je resultaten? Wat is de relevantie voor het werkveld?
%
% Daarom bestaat een abstract uit volgende componenten:
%
% - inleiding + kaderen thema
% - probleemstelling
% - (centrale) onderzoeksvraag
% - onderzoeksdoelstelling
% - methodologie
% - resultaten (beperk tot de belangrijkste, relevant voor de onderzoeksvraag)
% - conclusies, aanbevelingen, beperkingen
%
% LET OP! Een samenvatting is GEEN voorwoord!

%%---------- Nederlandse samenvatting -----------------------------------------
%
% Als je je bachelorproef in het Engels schrijft, moet je eerst een
% Nederlandse samenvatting invoegen. Haal daarvoor onderstaande code uit
% commentaar.
% Wie zijn bachelorproef in het Nederlands schrijft, kan dit negeren, de inhoud
% wordt niet in het document ingevoegd.

\IfLanguageName{english}{%
\selectlanguage{dutch}
\chapter*{Samenvatting}
\lipsum[1-4]
\selectlanguage{english}
}{}

%%---------- Samenvatting -----------------------------------------------------
% De samenvatting in de hoofdtaal van het document

\chapter*{\IfLanguageName{dutch}{Samenvatting}{Abstract}}

% 1. Waarover gaat de bachelorproef?
Deze bachelorproef gaat over automatisatie op het mainframe, meer specifiek een vergelijking en analyse van Rexx en Python. Rexx bestaat al heel lang op het mainframe, terwijl Python in vergelijking een nieuwkomer is. Door de skills gap en door de ontwikkeling van nieuwe technologieën buiten het mainframe, wordt er echter geargumenteerd dat Python Rexx kan vervangen als meest gebruikte automatisatie taal op het mainframe.
% 2. Waarom heb je er over geschreven?
Het doel van dit onderzoek is achterhalen hoe Rexx en Python zich qua automatisatie op z/OS verhouden vanuit technisch en praktisch oogpunt. Daarvoor wordt onderzocht wat automatisatie precies doet op het mainframe en wat de voornaamste redenen zijn om het mainframe te moderniseren. Daarna wordt een vergelijking van Python en Rexx gemaakt op basis van vooropgestelde use cases en requirements. Uiteindelijk wordt bepaald of Python een waardige vervanger kan zijn van Rexx.
% 3. Hoe heb je het onderzoek uitgevoerd?
Dit onderzoek is uitgevoerd aan de hand van een literatuurstudie, een requirementsanalyse en verschillende proof-of-concepts. Voor de requirementsanalyse zijn enkele interviews uitgevoerd en is een enquête opgesteld. In de proof-of-concepts worden use cases uitgevoerd met criteria bepaald aan de hand van de resultaten van de enquête en in samenspraak met de co-promotor van dit onderzoek.
% 4. Wat waren de resultaten? Wat blijkt uit je onderzoek?
Uit het onderzoek blijkt dat Rexx en Python elk hun eigen kwaliteiten hebben. Python wordt steeds populairder, en wordt meer en meer ondersteund op het mainframe, waardoor verschillende use cases van Rexx kunnen worden overgenomen. Rexx is echter zo verweven met de tools op het mainframe dat het (voorlopig) nog steeds onmisbaar is. Voor use cases dicht tegen het mainframe systeem geniet Rexx de voorkeur. Voor use cases waarbij communicatie naar de buitenwereld nodig is, is Python duidelijk beter. Op basis van dit onderzoek is de aanbeveling voor bedrijven om te moderniseren waar mogelijk, om op die manier nieuw talent aan te trekken. Dat talent kan nadien worden getraind met een mainframe achtergrond, inclusief Rexx, en kan op die manier bijdragen aan het onderhoud van de bestaande, en de ontwikkeling van de nieuwe, zaken op het mainframe.
Door enkele beperkingen van het gebruikte systeem en de tijd voor dit onderzoek zijn niet alle mogelijke use cases of requirements gebruikt om de talen met elkaar te vergelijken. Met meer tijd en een beter systeem zouden veel meer use cases kunnen worden uitgevoerd om een meer gedetailleerde vergelijking te doen. Daarnaast zou in verder onderzoek ook ruimer naar automatisatie kunnen worden gekeken door andere talen en technologieën te betrekken in de vergelijking. 
% 5. Wat betekenen je resultaten? Wat is de relevantie voor het werkveld?

% Daarom bestaat een abstract uit volgende componenten:
%
% - inleiding + kaderen thema
% - probleemstelling
% - (centrale) onderzoeksvraag
% - onderzoeksdoelstelling
% - methodologie
% - resultaten (beperk tot de belangrijkste, relevant voor de onderzoeksvraag)
% - conclusies, aanbevelingen, beperkingen
